\section{Introduction}\label{sec:intro}

The diagnosis of epilepsy rests on clinical history and electroencephalography (EEG), but the diagnostic value of a single routine EEG is limited.
The sensitivity of routine EEG after a first unprovoked seizure has been reported at only 17\% in adults \cite{lemoine2024}, only a small proportion of routine EEGs show interictal epileptiform discharges, and misdiagnosis of epilepsy has been reported at 20\% to 30\% \cite{myers2025}.
Computer analysis of routine EEG has been proposed to find markers that readers cannot see, but a systematic review of 37 studies concluded that none has entered clinical practice and that only 6 studies (16\%) used at least 100 subjects \cite{lemoine2024}.
For health informatics, the question is whether performance figures computed on recordings already stored in clinical archives can be trusted before model outputs are connected to clinical information workflows.
The intended readers are researchers and managers of clinical EEG archives who evaluate such figures, not clinicians making decisions about patients.

The public \textit{Temple University Hospital} (TUH) corpus and its epilepsy subset (TUEP) provide heterogeneous clinical EEG that is available to researchers for evaluation \cite{obeid2016,veloso2017}.
Studies on TUEP have reported accuracies above 95\% with a random forest on power spectra, using a random 67/33 split with no stated patient-level separation \cite{rohira2023}, and with a recurrent network on trajectory signatures, using five-fold cross-validation in which each segment is tested once \cite{tang2024}.
Such figures are difficult to interpret as disease information when segments from the same patient are in both the training and the test data: in EEG of patients with Alzheimer's disease, an accuracy of 99.8\% under a segment-wise split fell to 53.0\% under a subject-wise split \cite{brookshire2024}, and per-sample validation on a pathology task reached a balanced accuracy of almost 100\% where subject-wise schemes gave much lower values \cite{delpup2025}.
Routine clinical data also contain acquisition biases that correlate with the label, which models learn as shortcuts and which internal cross-validation does not detect \cite{ongly2024}. In TUEP, the epilepsy group is younger and was recorded more often in particular eras and on particular devices (Subsection~\ref{sec:populasi}). TUEP is also not a first-seizure population: patients with epilepsy have a median of three sessions in the corpus, and the label takes medication at the time of recording into account.

A common response is to remove these covariates from the features, for example by regression or normative modeling \cite{rutherford2022,rutherford2023,bayer2022neuroimage}.
Removing confounders by per-feature regression can, however, move confounder information into non-normal residuals that nonlinear classifiers exploit, so that a null or moderate result becomes almost perfect; Hamdan et al. showed this on tabular data and proposed to shuffle the features before removal to reveal it \cite{hamdan2023}. Statistical tests for confounding bias in predictive models are also available \cite{spisak2022}.
Clinical EEG differs from that setting in three ways that matter for an audit: each patient contributes many windows, so the data are clustered; the test produces a distribution of AUCs over replicates that has to be turned into a decision, and in the references reviewed the error frequencies of such a decision have not been reported; and the representations range from patient-level spectra to sequence models, for which the mitigations have not been compared there.
When a method is then chosen after many rounds on the same patients, selection optimism adds to these problems \cite{varma2006}.

This study asks one primary question: when the EEG carries no information about the label, how much discrimination does covariate adjustment create in patient-level EEG classification, and how reliably do explicit decision rules detect it?
The main effect size is the paired difference in AUC between the same pipeline with and without adjustment under a null experiment (in percentage points); \tabref{tab:plasmode} answers it on real EEG features and \tabref{tab:oc} gives the error frequencies of the rules.
The case study is TUEP version 2.0.1 (200 patients, first sessions only, the same 50 patient-level splits for all case-study models), and the object of the audit is a three-view fusion pipeline with covariate adjustment (F2), whose originally planned adjustment turned out to leak.
Secondary aims were to place the audited pipeline against non-EEG baselines and within matched strata, to measure the EEG information beyond all non-EEG predictors, and to estimate selection optimism conditional on the stored candidates. The study does not claim a better model or clinical value.
\tabref{tab:matriks} separates what was known from what this study adds.

\input{\naskahdir/tabel_matriks_kontribusi}
